\documentclass[11pt,a4paper]{article}

\usepackage[margin=2.5cm]{geometry}
\usepackage{amsmath,amssymb,amsthm}
\usepackage{booktabs}

\newtheorem{proposition}{Proposition}[section]

\newcommand{\R}{\mathbb{R}}
\newcommand{\Z}{\mathbb{Z}}
\newcommand{\Q}{\mathbb{Q}}

\title{\textbf{The Dual Meixner--Pollaczek Matching Coefficients:\\
the Calculation}}
\author{Robin Langer and Claude Code}
\date{August 2026}

\begin{document}
\maketitle

%% ===================================================================
\section{The Problem}
%% ===================================================================

Let $P_n$ be the Meixner--Pollaczek polynomials
($\lambda = \tfrac14$, $\phi = \tfrac{\pi}{2}$),
\[
  (n+1) P_{n+1} = 2x P_n - (n - \tfrac12) P_{n-1},
  \qquad P_0 = 1, \quad P_1 = 2x,
\]
orthogonal for the weight
$W(x) = \Gamma(\tfrac14+ix)\,\Gamma(\tfrac14-ix)$ on $\R$, with
$\int P_m P_k W = h_k \delta_{mk}$, $h_0 = \sqrt2\,\pi^{3/2}$, and
$\widehat{P}_n = P_n/\sqrt{h_n}$.

The weight factors into an analytic function and its reflection, and
distributing one factor to each side defines a \emph{pair of dual
Meixner--Pollaczek functions}
\[
  F_m(x) = \Gamma(\tfrac14+ix)\, \widehat{P}_m(x),
  \qquad
  G_m(x) = \Gamma(\tfrac14-ix)\, \widehat{P}_m(x),
\]
biorthonormal under the bilinear pairing:
$\int_\R F_m G_k\, dx = \delta_{mk}$.

\textbf{Problem.} Approximate each dual function by
Meixner--Pollaczek polynomials: find the coefficients in the best
$L^2(W)$-approximation
\[
  F_m \approx \sum_{k \le N} c_{mk}\, \widehat{P}_k,
  \qquad
  c_{mk} = \langle F_m, \widehat{P}_k \rangle_W
  = \frac{J_{mk}}{\sqrt{h_m h_k}},
\]
where
\begin{equation}\label{eq:J}
  J_{mk} = \int_\R P_m(x)\, P_k(x)\;
  \Gamma(\tfrac14+ix)^2\, \Gamma(\tfrac14-ix)\; dx.
\end{equation}
The coefficients for $G_m$ are the complex conjugates, so it suffices
to compute $J_{mk}$. We do so in closed form.

%% ===================================================================
\section{The Three-Gamma Integral}
%% ===================================================================

\begin{proposition}\label{prop:3gamma}
For $\operatorname{Re} a, b, c > 0$,
\[
  \int_\R \Gamma(a+ix)\,\Gamma(b+ix)\,\Gamma(c-ix)\, dx
  = 2\pi\, \Gamma(a+c)\,\Gamma(b+c)\; U(a+c,\, a-b+1,\, 1),
\]
where $U$ is Tricomi's confluent hypergeometric function.
\end{proposition}

\begin{proof}
Write each factor as $\Gamma(s) = \int_0^\infty t^{s-1} e^{-t}\, dt$
with variables $t, u$ (for $+ix$) and $v$ (for $-ix$). The $x$-integral
gives $\int_\R (tu/v)^{ix} dx = 2\pi\,\delta(\log(tu/v))$, which
collapses the $v$-integral at $v = tu$:
\[
  \int_\R
  = 2\pi \int_0^\infty\!\!\int_0^\infty
    t^{a+c-1}\, u^{b+c-1}\, e^{-t-u-tu}\, dt\, du.
\]
The $u$-integral is $\Gamma(b+c)(1+t)^{-(b+c)}$, and the remaining
$t$-integral is Tricomi's representation
$\Gamma(\alpha)\, U(\alpha, \beta, 1)
= \int_0^\infty t^{\alpha-1}(1+t)^{\beta-\alpha-1} e^{-t}\, dt$
with $\alpha = a+c$, $\beta = a-b+1$.
\end{proof}

The same proof with an oscillatory factor $e^{isx}$ inserts $e^{-s}$
into the delta constraint ($v = tu e^{-s}$) and gives the shifted form
\begin{equation}\label{eq:shifted}
  \int_\R \Gamma(a+ix)\,\Gamma(b+ix)\,\Gamma(c-ix)\, e^{isx}\, dx
  = 2\pi\, e^{-as}\, \Gamma(a+c)\,\Gamma(b+c)\;
    U(a+c,\, a-b+1,\, e^{-s}).
\end{equation}

At $a = b = c = \tfrac14$, using
$U(\tfrac12, 1, 2z) = e^{z} K_0(z)/\sqrt{\pi}$,
\begin{equation}\label{eq:base}
  J_{00} = 2\pi\, \Gamma(\tfrac12)^2\, U(\tfrac12, 1, 1)
  = 2\pi^{3/2} \sqrt{e}\; K_0(\tfrac12)
  = 16.97348192607767\ldots
\end{equation}
(verified to 20 digits by quadrature).

%% ===================================================================
\section{Reduction of the Polynomial Moments}
%% ===================================================================

The recurrence $(a+ix)\,\Gamma(a+ix) = \Gamma(a+1+ix)$ converts
multiplication by $x$ into integer shifts of the parameters:
\[
  ix\,\Gamma(a+ix) = \Gamma(a+1+ix) - a\,\Gamma(a+ix),
  \qquad
  -ix\,\Gamma(c-ix) = \Gamma(c+1-ix) - c\,\Gamma(c-ix).
\]
Applying these repeatedly, every moment
$\int x^j\, \Gamma(\tfrac14+ix)^2 \Gamma(\tfrac14-ix)\, dx$ is a
$\Q$-combination of integrals of Proposition~\ref{prop:3gamma} with parameters in
$\tfrac14 + \mathbb{N}$, hence of values $U(\alpha, \beta, 1)$ with
$\alpha \in \tfrac12 + \mathbb{N}$, $\beta \in \Z$. The contiguous
relations for $U$ reduce all of these to the $\Q$-span of two
generators, which we take (via
$K_\nu(z) = \sqrt{\pi}\,(2z)^\nu e^{-z}\, U(\nu+\tfrac12, 2\nu+1, 2z)$)
to be
\[
  A = \sqrt{2e}\; K_0(\tfrac12) = 2.155416163\ldots, \qquad
  B = \sqrt{2e}\; K_1(\tfrac12) = 3.862230968\ldots
\]
Since the $P_m$ have rational coefficients, every $J_{mk}/h_0$ lies in
$\Q A \oplus \Q B$. Empirically the denominators are exactly
$2^m m!\; 2^k k!$:
\begin{equation}\label{eq:T}
  T_{mk} = \frac{2^m m!\, 2^k k!}{h_0}\, J_{mk}
  \;\in\; \Z A \oplus \Z B
  \qquad (\text{$T_{mk}$ real for $m+k$ even, } i \cdot \text{real
  for $m+k$ odd}).
\end{equation}

%% ===================================================================
\section{The Generating Function}
%% ===================================================================

The whole table has a closed generating function. Substituting the
Meixner--Pollaczek generating function
\[
  \sum_{n \ge 0} P_n(x)\, w^n
  = (1-iw)^{-\frac14+ix} (1+iw)^{-\frac14-ix}
\]
once in each index of~\eqref{eq:J} turns all $x$-dependence into a
single oscillation $e^{isx}$ with
\[
  e^{-s} = \zeta(w,u) = \varphi(w)\,\varphi(u),
  \qquad
  \varphi(t) = \frac{1+it}{1-it}
  \quad (\text{the Cayley transform}),
\]
and~\eqref{eq:shifted} evaluates the integral. The algebraic
prefactors collapse, and using
$U(\tfrac12, 1, 2z) = e^{z} K_0(z)/\sqrt{\pi}$ once more,
\begin{equation}\label{eq:gf}
  \sum_{m,k \ge 0} J_{mk}\, w^m u^k
  = \frac{2\pi^{3/2}\, e^{\zeta/2}\, K_0(\zeta/2)}
         {\sqrt{(1-iw)(1-iu)}}
\end{equation}
(verified numerically to $15$ digits against every entry of the
table). At $w = u = 0$, $\zeta = 1$, recovering~\eqref{eq:base}.

Formula~\eqref{eq:gf} explains the two-generator structure: every
Taylor coefficient acts on $K_0$ through the closed derivative algebra
$K_0' = -K_1$, $K_1' = -K_0 - K_1/z$, evaluated at $z = \tfrac12$, so
each $J_{mk}/h_0$ lies in $\Q A \oplus \Q B$ automatically (proving
what \S3 stated empirically, up to the denominators). It also makes
the real/imaginary checkerboard of~\eqref{eq:T} visible by inspection:
$(w,u) \mapsto (-w,-u)$ conjugates $\zeta$.

The derivative algebra is a weighted \emph{two-state automaton} on
states $\{K_0, K_1\}$: each differentiation either switches state
(weight $-1$) or loops at $K_1$ (weight $-1/z$, i.e.\ $-2$ at
$z = \tfrac12$). Expanding~\eqref{eq:gf} by Fa\`a di Bruno, the
integer coefficients of $A$ and $B$ in $T_{mk}$ are thus signed sums
over pairs (set partition of the $m+k$ derivatives, path in the
automaton with one step per block), each derivative dressed by a
Taylor coefficient of $\zeta = \varphi(w)\varphi(u)$ at the Cayley
transform --- the source of the powers of $2$ absorbed by the
normalisation $2^m m!\, 2^k k!$ in~\eqref{eq:T}. A self-contained
combinatorial formula for the pair of integer matrices remains to be
extracted.

%% ===================================================================
\section{The Coefficients}
%% ===================================================================

The table lists the rescaled integrals $T_{mk}$ of~\eqref{eq:T}, with
integer coefficients recovered by PSLQ at $40$ digits (the factor $i$
for $m+k$ odd is dropped); $(3,5)$ exceeded the coefficient search
bound. To recover the approximation coefficients of $F_m$: since
$h_n = h_0 \binom{2n}{n}/4^n$, the normalisations collapse to
\[
  c_{mk}
  = \frac{T_{mk}}{m!\, k!\,
    \sqrt{\binom{2m}{m}\binom{2k}{k}}}\,,
\]
and the coefficients of $G_m$ are the complex conjugates.

\begin{table}[ht]
\centering
\small
\begin{tabular}{cl@{\qquad}cl}
\toprule
$(m,k)$ & $T_{mk}$ & $(m,k)$ & $T_{mk}$ \\
\midrule
$(0,0)$ & $A$                   & $(1,3)$ & $2575A - 1440B$ \\
$(0,1)$ & $3A - 2B$             & $(1,4)$ & $39411A - 21978B$ \\
$(0,2)$ & $12B - 23A$           & $(1,5)$ & $394932B - 707389A$ \\
$(0,3)$ & $138B - 241A$         & $(2,2)$ & $2369A - 1320B$ \\
$(0,4)$ & $3265A - 1800B$       & $(2,3)$ & $33743A - 18834B$ \\
$(0,5)$ & $52907A - 29706B$     & $(2,4)$ & $317892B - 569719A$ \\
$(1,1)$ & $12B - 21A$           & $(2,5)$ & $6190602B - 11092229A$ \\
$(1,2)$ & $114B - 205A$         & $(3,3)$ & $295572B - 529549A$ \\
        &                       & $(3,4)$ & $5399634B - 9675553A$ \\
\bottomrule
\end{tabular}
\end{table}

%% ===================================================================
\section{The Reverse Expansion}
%% ===================================================================

Expanding in the other direction — the polynomials in terms of the
dual functions — is elementary by contrast. Biorthogonality gives
\[
  \widehat{P}_m = \sum_k D_{mk} F_k,
  \qquad
  D_{mk} = \int_\R \widehat{P}_m\, G_k\, dx
  = \frac{1}{\sqrt{h_m h_k}}
    \int_\R P_m(x)\, P_k(x)\, \Gamma(\tfrac14-ix)\, dx,
\]
a \emph{one}-Gamma integral. Mellin inversion of
$\Gamma(s) = \int_0^\infty t^{s-1} e^{-t} dt$ gives
$\int_\R \Gamma(a+ix)\, t^{-ix} dx = 2\pi\, t^a e^{-t}$, and the same
substitution of the generating function as in \S4 yields
\begin{equation}\label{eq:gfrev}
  \sum_{m,k \ge 0}
  \Big( \textstyle\int_\R P_m P_k\, \Gamma(\tfrac14-ix)\, dx \Big)
  w^m u^k
  = \frac{2\pi\, e^{-1/\zeta}}{\sqrt{(1+iw)(1+iu)}}
\end{equation}
(verified to $20$ digits; at $w=u=0$ this is $2\pi/e$). Where the
forward direction needed the Bessel pair $(A, B)$, the reverse needs
the single generator $e^{-1}$: one Gamma factor inverts to a pure
exponential, whose derivative algebra has one state, while three
Gamma factors produce the Tricomi $U$, whose Bessel algebra has two.
The two generating functions are exchanged by
$\zeta \mapsto 1/\zeta$, $(w,u) \mapsto (-w,-u)$ --- the Cayley
inversion swapping the half-planes.

The same normalisation as in~\eqref{eq:T} again produces integers ---
now over the \emph{single} generator $2\pi/e$, and provably so, since
the Taylor coefficients of~\eqref{eq:gfrev} are Gaussian rationals
times $e^{-1}$ and can be expanded symbolically (no PSLQ needed):
\[
  S_{mk} = 2^m m!\, 2^k k!\; \frac{e}{2\pi}
  \int_\R P_m P_k\, \Gamma(\tfrac14 - ix)\, dx \;\in\; \Z,
  \qquad
  D_{mk} = \frac{1}{e}\sqrt{\frac{2}{\pi}}\;
  \frac{S_{mk}}{m!\, k!\, \sqrt{\binom{2m}{m}\binom{2k}{k}}}
\]
(again dropping a factor $i$ when $m+k$ is odd).

\begin{table}[ht]
\centering
\small
\begin{tabular}{r|rrrrrr}
$S_{mk}$ & $k=0$ & $1$ & $2$ & $3$ & $4$ & $5$ \\
\hline
$m=0$ & $1$ & $3$ & $5$ & $11$ & $249$ & $-2621$ \\
$1$ &  & $7$ & $47$ & $399$ & $-2005$ & $263$ \\
$2$ &  &  & $473$ & $-617$ & $22109$ & $-402065$ \\
$3$ &  &  &  & $31511$ & $-412557$ & $-3017969$ \\
$4$ &  &  &  &  & $-1916623$ & $-29235989$ \\
$5$ &  &  &  &  &  & $-1012892921$ \\
\end{tabular}
\end{table}

These integers have a combinatorial interpretation. In the rescaled
variables $v = 2iw$, $\mu = 2iu$ the generating
function~\eqref{eq:gfrev} factors as
\[
  \bigl(\sigma'(v)\,\tau'(\mu)\bigr)^{1/4}\, e^{\sigma+\tau-\sigma\tau},
  \qquad
  \sigma = \frac{v}{1+v/2}, \quad \tau = \frac{\mu}{1+\mu/2}
\]
--- a Charlier kernel in Laguerre variables, with the
Meixner--Pollaczek parameter $\lambda = \tfrac14$ appearing as the
exponent of the Jacobian prefactor. Reading this as a species:

\begin{proposition}
$S_{mk} = (-1)^{\lfloor (m+k)/2 \rfloor} \sum_{\text{structures}}
\operatorname{wt}$, summed over structures on $([m], [k])$ in which
each of $[m]$, $[k]$ is split into a \emph{defect set} and a forest of
nonempty \emph{chains} (linear orders), together with a partial
matching between the $[m]$-chains and the $[k]$-chains; the weights
are $(-1)^d (2d-1)!!$ for a defect set of size $d$ (the prefactor),
$(-1)^{\ell-1}\, 2^{\ell+1}$ for a chain of length $\ell$ (the
substitution), and $-1$ per matched pair of chains (the Charlier
kernel $e^{\sigma+\tau-\sigma\tau}$).
\end{proposition}

Equivalently, with Lah numbers
$L(n,c) = \binom{n-1}{c-1} n!/c!$ counting the chain forests,
\[
  S_{mk} = (-1)^{\lfloor \frac{m+k}{2} \rfloor}
  \sum_{d_1, c_1, d_2, c_2}
  W_m(d_1, c_1)\, W_k(d_2, c_2)\; C_{c_1 c_2},
\]
\[
  W_n(d, c) = \binom{n}{d} (-1)^{d} (2d-1)!!\;
  L(n{-}d,\, c)\, (-1)^{n-d-c}\, 2^{n-d+c},
\]
where $C_{c_1 c_2} = \sum_j (-1)^j \binom{c_1}{j}\binom{c_2}{j} j!$
is the signed count of partial matchings (a Charlier polynomial
value). This reproduces the entire table exactly.

One caution: the coefficients $D_{mk}$ \emph{grow} along the diagonal
($D_{00} = 0.2935\ldots$, $D_{11} = 1.027\ldots$,
$D_{22} = 5.785\ldots$), so the reverse expansion is formal —
$\widehat{P}_m$ does not lie in the closed span of the $F_k$. The
analytic content lives in the forward direction; the reverse carries
the simpler arithmetic.

\end{document}
